\section{Constructing the signed reference difference}
\label{sec:reference-difference}
The cone in Section~\ref{sec:reference-cone} applies only to a
comparison produced on a common background history. We describe
that construction. Its principal feature is to remove the
coefficient vector after each backward step, retaining the removed
part in a reference label. This provides a family of exact
comparisons, from which a single label is chosen at the coarsest
scale.

Fix the degree $d$ and scales $l\le k$. Assume the bank and
column inequalities \eqref{eq:reference-bank}, with $K_i,C_i\ge1$,
as in Section~\ref{sec:reference-cone}. At each scale $i$, let
$\mathcal P_i(b)$ be the calibrated reference coefficient field
with bare label $b$ in the fixed coefficient space. Its coefficient
vector is $I_i b$, where $I_i$ is a linear isomorphism. Assume its
density, ordinary field and flagged profile depend linearly on
$b$, and its remainder has a nonnegative subadditive bound $h_i(b)$
with
\begin{equation}\label{eq:reference-primary-charge}
           h_i(b)\le C_i\mu(\operatorname{core}I_i b)
                         \le C_i\mu(I_i b).
\end{equation}
The reference densities satisfy the exact degree-$d$ transition
identity. Their coefficient and remainder data satisfy
\eqref{eq:reference-direction}, with a common comparison map for
each reference label, as supplied by the calibrated construction.

We also require a small transition estimate for a coefficient
field $S$ whose coefficient vector is zero. If its remainder is
bounded by $q\ge0$, its backward transition has a representation
$U$ such that
\begin{equation}\label{eq:reference-zero-transition}
 \operatorname{dens}(U)=L_i^d\mathcal T_i\operatorname{dens}(S),
 \quad \mu(U.x)\le\eta_i q,\quad
 \operatorname{rem}(U)\le\eta_i q,
 \quad (1+2C_i)\eta_i\le\varepsilon_i.
\end{equation}
Here $\eta_i\ge0$ and $\varepsilon_i=L_i^{-3/8}$.
The fields and their sums retain the coefficient-field format and
the integrability needed for the displayed transition identities.
These are analytic inputs for the same history, including its
calibrated reference fields.

\subsection{Removing the coefficient after each step}
Start with a coefficient field $Q$ at scale $k$, with remainder
bound $t\ge0$ and $\mu(Q.x)+t\le w$. Put
\[
 c_k=I_k^{-1}Q.x,\qquad
 S_k=Q-\mathcal P_k(c_k),\qquad
 q_k=t+h_k(-c_k).
\]
The minus sign denotes a sum with the field $\mathcal P_k(-c_k)$;
its regular remainders are concatenated with those of $Q$.
Thus $S_k.x=0$ and its remainder is at most $q_k$. For every
bare label $v$, define
\[
 D_k(v)=S_k+\mathcal P_k(c_k+v),\qquad
 r_k(v)=q_k+h_k(c_k+v).
\]
Linearity of the reference density and vector gives
\begin{equation}\label{eq:reference-top-family}
 D_k(v).x=Q.x+I_kv,
 \qquad\operatorname{dens}(D_k(v))
                  =\operatorname{dens}(Q)+\operatorname{dens}(\mathcal P_k(v)).
\end{equation}
The cost, despite possible cancellation in this identity, is
kept as the sum of the displayed nonnegative remainder bounds.
Since $t\le w$ and $\mu(Q.x)\le w$,
\[
 r_k(v)\le w+C_k\{w+\mu(D_k(v).x)\}.
\]
This is exactly the finest-scale condition
\eqref{eq:reference-terminal}.

Suppose $c_{i+1},S_{i+1},q_{i+1}$ have been constructed, with
$S_{i+1}.x=0$. Apply \eqref{eq:reference-zero-transition} to obtain
$U_i$, put $e_i=\eta_iq_{i+1}$, and set
\begin{align*}
 b_i&=I_i^{-1}U_i.x,&c_i&=c_{i+1}+b_i,\\
 S_i&=U_i-\mathcal P_i(b_i),&
 q_i&=e_i+h_i(-b_i).
\end{align*}
Then $S_i.x=0$, and $q_i$ bounds its remainder. Define
$D_i(v)=S_i+\mathcal P_i(c_i+v)$ and
$r_i(v)=q_i+h_i(c_i+v)$. Their densities satisfy
\[
 \operatorname{dens}(D_i(v))
           =\operatorname{dens}(U_i)
                       +\operatorname{dens}(\mathcal P_i(c_{i+1}+v))
           =L_i^d\mathcal T_i\operatorname{dens}(D_{i+1}(v)).
\]
The first equality cancels the two reference terms with labels
$-b_i$ and $b_i$ by linearity. The second uses both transition
identities and linearity of the integrable transition. The same
cancellation holds separately for ordinary fields and flagged
profiles; it is not merely an equality after taking expectation.

We verify the signed one-step estimates. Fix $v$ and write
$y=c_{i+1}+v$. Then
\[
 D_{i+1}(v).x=I_{i+1}y,\qquad D_i(v).x=I_i y+U_i.x.
\]
Use the comparison map belonging to the primary label $y$.
Its coefficient error is at most
\[
 \varepsilon_i\{\mu(\operatorname{core}I_{i+1}y)+h_{i+1}(y)\}.
\]
The added vector $U_i.x$ costs at most
$\eta_iq_{i+1}\le\varepsilon_iq_{i+1}$. Thus the coefficient
inequality in \eqref{eq:reference-direction} holds for the
actual signed vectors, with cost $r_{i+1}(v)$.

For the remainder, \eqref{eq:reference-primary-charge} gives
$h_i(\pm b_i)\le C_i e_i$. Subadditivity yields
\begin{align*}
 r_i(v)&=e_i+h_i(-b_i)+h_i(y+b_i)\\
       &\le h_i(y)+(1+2C_i)e_i\\
       &\le K_i\mu(\operatorname{core}D_{L_i}I_{i+1}y)
                         +\varepsilon_i\{h_{i+1}(y)+q_{i+1}\}.
\end{align*}
This is the other inequality of
\eqref{eq:reference-direction}. In particular it is a bound
for the very same presentation that supplies the exact density
identity. No subtraction of unsigned remainder bounds is used.

Now choose the single label $v=-c_l$ and write $b=c_l$.
The coefficient vector of $D_l(-b)$ vanishes. The construction
therefore meets every hypothesis of
Section~\ref{sec:reference-cone}, including its terminal condition.
Writing $Z_i=D_i(-b)$ and $r_i=r_i(-b)$ gives
\begin{equation}\label{eq:reference-produced-difference}
 \operatorname{dens}(Z_k)
       =\operatorname{dens}(Q)-\operatorname{dens}(\mathcal P_k(b)),
 \qquad
 r_i+\mu(Z_i.x)\le\frac{4(1+C_k)w}{e^{(s_k-s_i)/8}}.
\end{equation}
If $E\ge4(1+C_k)w$, then $w\le E$ and at $i=k$
\[
              \mu(I_k b)=\mu(Q.x-Z_k.x)\le w+E\le2E.
\]
This is the quota bound for the matched reference label.

\subsection{The singleton transition expansion is retained}
To use this difference inside a marked density, more than its
coefficient and remainder estimates is needed. Suppose the
zero-vector transition in \eqref{eq:reference-zero-transition}
comes with the singleton expansion of
Section~\ref{sec:local-transfer}, of fine cost $2q_{i+1}$ and
coarse response cost $e_i$. Suppose the primary transition for
label $y$ has fine cost at most
\[
 \operatorname{orb}(d)\mu(I_{i+1}y)
                            (\psi_i+2C_{i+1}),
\]
and coarse response cost $h_i(y)$. Both expansions must use the
same quoted-profile map; their finite sets of slots can then be
united. Their sum has the ordinary and flagged profiles of
$D_{i+1}(v)$ and the coarse profile of $D_i(v)$, by the
cancellations just proved. Its fine cost is at most
\[
 \operatorname{orb}(d)\mu(D_{i+1}(v).x)
           (\psi_i+2C_{i+1})+2r_{i+1}(v).
\]
Its response cost $e_i+h_i(y)$ is at most $r_i(v)$, because
\[
 h_i(y)=h_i\bigl(-b_i+(y+b_i)\bigr)
                           \le h_i(-b_i)+h_i(y+b_i).
\]
Thus the same signed difference carries the singleton transition
expansions needed for the marked transport argument. Combining
\eqref{eq:reference-produced-difference} with
\eqref{eq:reference-affordable} makes these expansions affordable
once the prescribed one-step allowances are imposed. The
production of the small zero-vector transition and the calibrated
primary transition bounds remains an analytic dependency of this
construction.
